\section{Fazit}
Insgesamt lieferten beide Teile des Verusuchs Ergebnisse, die mit der Hypothese übereinstimmten. Zunächst wurde die spektrale Verteilung der kosmischen Strahlung bestimmt, wobei eine Korrelation in $cos^2(\theta)$ festgestellt wurde. Der Exponent betrug:
\begin{align*}
    c=(\pm)
\end{align*}

Diskussion
(Die Anzahl der zufälligen Koinzidenzen war ebenfalls konsistent:) Es wurden mehr/weniger geometrische Koinzidenzen (???) als zufällige zeitliche Koinzidenzen (???) registriert, was auf eine korrekte Nutzung der Koinzidenzlogik hindeutet.

Das aufgezeichnete Pulshöhenspektrum entspricht eine Landau-Verteilung, ähnlich der von dünnen Szintillatoren. Die Landau-Anpassung ergab:
\begin{align*}
    \chi^2_{net} = ????
\end{align*}
( besser als bei der Gauß-Verteilung mit dem Wert $\chi^2_{net} = ????$, \\
mehr Diskussion\\ 
Im zweiten Teil des Versuchs wird die Lebensdauer der Myonen bestimmt. Der erhaltene Wert
$\tau=(\pm)\mu s$ \\
mehr Diskussion\\
Die während des Veruchs identifizierten systematischen Unsicherheiten und Fehlerquellen, insbesondere jene, die mit der Triggerlogik, Hintergrundereignissen und statistischen Daten zusammenhängen; erklären wahrscheinlich diese Diskrepanz. Dennoch kann betrachtet werden, da das fundamentale Zerfallsgesetz experimentell reproduziert und die Messmethode validiert wurden.